\section{Related work}\label{sec:related}

\textbf{CuTe}~\cite{cutedocs,cecka} is the system under study; its operations
and their conditions are stated in Section~\ref{sec:intro}.
\textbf{Shah}~\cite{shah} formalizes the algebra and gives the conditions
under which the composite of two layouts is again a layout. In our algebra
that is decided after the fact by Theorem~\ref{thm:decide}, and a composite
that is not a layout is still a legal map. \textbf{Carlisle et
al.}~\cite{categorical} model a class of CuTe layouts as morphisms of two
categories, define composition, logical product and logical division on the
morphisms, and characterize the layouts so obtained. \textbf{Linear
layouts}~\cite{triton}, in Triton, represent a layout as a matrix over
$\mathbb{F}_2$ acting on the bits of the hardware index, which gives generic
layout-to-layout conversion; a swizzle of Definition~\ref{def:swizzle} is such
a matrix. \textbf{Bhaskaracharya et al.}~\cite{isl} represent both CuTe
layouts and linear layouts as integer set relations in ISL and implement
composition, inversion and complement with its operations; the question of
Section~\ref{sec:decide} can be posed there, and Theorem~\ref{thm:decide}
answers it without a solver. \textbf{Axe}~\cite{axe} maps logical
coordinates to a multi-axis physical space through named axes and unifies
tiling, sharding, replication and offsets from device meshes to threads; its
replication is explicit, as $\broadcast$ is in our algebra, and its
Appendix~A.2.1 uses the floor expansion and first differences of
Section~\ref{sec:decide} to prove a normalized representation unique.
\textbf{Halide}~\cite{halide}, \textbf{TVM}~\cite{tvm} and
\textbf{MLIR}~\cite{mlir} carry general integer expression simplifiers;
Section~\ref{sec:decide} decides the layout question directly instead.
